\section{Experiments}
\label{sec:experiments}


We evaluate \ours along three dimensions: \emph{native policy performance},
\emph{component composability}, and \emph{local dynamics modeling}. We test
whether the shared robot-video foundation supports different video--action
interaction programs, whether independently trained action components remain
useful when paired with adapted visual predictors, and whether local-context forward
dynamics models action-dependent visual outcomes.


\subsection{Experimental Setup}
\label{sec:experimental-setup}

\paragraph{Robot-video pretraining and simulation.}
We initialize the visual backbone from Wan2.2-5B~\cite{wan2025wan} and further
pretrain it without action supervision on heterogeneous robot and
human-interaction video, including UMI/MV-UMI~\cite{chi2024universal,rayyan2026mv},
AgiBot World~\cite{bu2025agibot}, RoboMIND~\cite{wu2024robomind},
InternData-A1~\cite{tian2026interndata}, RoboCOIN~\cite{wu2025robocoin},
FastUMI~\cite{zhaxizhuoma2025fastumi}, Ego-Exo4D~\cite{grauman2024ego},
and Open X-Embodiment~\cite{o2024open}.
% Dataset composition, preprocessing, sampling, and optimization details are provided in the
% appendix.
Unless otherwise stated, downstream \ours variants use this
common initialization.


The four interaction programs share the same dual-stream
MoT backbone and downstream training recipe; they
differ in generation order and future cross-modal attention. Video is
processed in four-frame latent chunks, each predicting a 16-step action
chunk, with proprioception provided as a per-chunk conditioning signal.
Actions consist of end-effector position, axis-angle rotation, and a
gripper command. We use a two-stage training procedure: causal
robot-video pretraining followed by joint video--action training on the
target suite.


We evaluate Joint, VTA, ATV, and Decoupled on LIBERO-Object, LIBERO-Goal,
LIBERO-Spatial, and LIBERO-Long~\cite{liu2023libero}.
Component transfer uses LIBERO-Long as the downstream source
distribution and four LIBERO-90 targets (Table~\ref{tab:transfer-tasks}) chosen to represent qualitatively different forms of transfer. These targets are outside the LIBERO-Long
source-task set; robot-video pretraining is treated as a separate stage.





\begin{table}[t]
    \centering
    \caption{Real-world bimanual manipulation tasks.}
    % \vspace{-0.05in}
    \label{tab:real-world-tasks}
    \small
    \resizebox{\linewidth}{!}{%
    \begin{tabular}{@{}lll@{}}
        \toprule
        \textbf{Task}
        & \textbf{Objective}
        & \textbf{Main challenge} \\
        \midrule
        Toast Bread
        & Insert bread and activate toaster
        & Deformability and precise insertion \\
        Rubik's Cube
        & Restore final layer of $2\times2$ cube
        & Contact-rich push and rotation \\
        Sort Cups
        & Sort cups according to color
        & Variable ordering and sequencing \\
        \bottomrule
    \end{tabular}}
    % \vspace{-0.2in}
\end{table}

\paragraph{Real-world tasks.}
We evaluate on a bimanual platform comprising two Franka Research 3 arms
with parallel-jaw grippers and three RGB views: two wrist-mounted cameras
and one centered third-person camera. The three task definitions
are listed in Table~\ref{tab:real-world-tasks}. Figure~\ref{fig:tasks} shows the
real-world tasks and the LIBERO-90 transfer targets.

\paragraph{Dynamics supervision.}
Native task models are trained from task demonstrations. For local dynamics,
we additionally construct \textbf{LIBERO-Long-CF} by executing alternative
action continuations from states encountered in LIBERO-Long demonstrations.
Branches sharing an initial state remain in the same train, validation, or
test partition.
IDM and FDM experiments use demonstration-only, counterfactual-only
(CF-only), or mixed supervision. Mixed supervision samples 60\% of
transitions from $\mathcal D_{\mathrm{cf}}$ and 40\% from
$\mathcal D_{\mathrm{task}}^{\mathrm{loc}}$, corresponding to $\eta=0.6$
in Eq.~\ref{eq:dynamics-training-distribution}.

%For physical experiments, the broader dynamics data instead combines successful demonstrations, failed attempts, and free-form play.
% Collection and sampling details are provided in the appendix.

\paragraph{Evaluation.}
Closed-loop task success is the primary policy and composition metric.
For forward dynamics, we report visual prediction error and outcome
identification accuracy. Given $K$ realized futures from the same state,
the metric tests whether the prediction is closest in RGB MSE to its
matching outcome. We report both pairwise identification ($K=2$) and
identification among all 16 action branches ($K=16$).


\begin{table}[!t]
    \centering
    \caption{Closed-loop success (\%) on four LIBERO suites.
Baselines are previously reported results.}
    % \vspace{-0.1in}
    \label{tab:libero-broad}
    \scriptsize
    \setlength{\tabcolsep}{5pt}
    \begin{tabular}{@{}lccccc@{}}
        \toprule
        \textbf{Method}
        & \textbf{Object}
        & \textbf{Goal}
        & \textbf{Spatial}
        & \textbf{Long}
        & \textbf{Mean} \\
        \midrule
        OpenVLA~\cite{kim2024openvla}
        & 88.4 & 79.2 & 84.7
        & 53.7 & 76.5 \\
        OpenVLA-OFT~\cite{kim2025fine}
        & 98.4 & 97.9 & 97.6
        & 94.5 & 97.1 \\
        $\pi_0$~\cite{black2024pi_0}
        & 98.8 & 95.8 & 96.8
        & 85.2 & 94.1 \\
        $\pi_{0.5}$~\cite{intelligence2025pi_}
        & 98.2 & 98.0 & 98.8
        & 92.4 & 96.9 \\
        GR00T-N1~\cite{bjorck2025gr00t}
        & 97.6 & 93.0 & 94.4
        & 90.6 & 93.9 \\
        Motus~\cite{bi2026motus}
        & 99.8 & 96.6 & 96.8
        & 97.6 & 97.7 \\
        Fast-WAM~\cite{yuan2026fast}
        & 100.0 & 97.0 & 98.2
        & 95.2 & 97.6 \\
        LingBot-VA~\cite{li2026causal}
        & 99.6 & 97.2 & 98.5
        & 98.5 & 98.5 \\
        \midrule
        \textbf{\ours-VTA}
        & $\text{99.4}_{\pm\text{0.3}}$
        & $\text{98.4}_{\pm\text{0.5}}$
        & $\text{98.6}_{\pm\text{0.2}}$
        & $\text{97.8}_{\pm\text{0.4}}$
        & 98.6 \\

        \textbf{\ours-ATV}
        & $\text{98.0}_{\pm\text{0.4}}$
        & $\text{97.2}_{\pm\text{0.2}}$
        & $\text{96.6}_{\pm\text{0.6}}$
        & $\text{95.4}_{\pm\text{0.3}}$
        & 96.8 \\

        \textbf{\ours-Joint}
        & $\text{98.2}_{\pm\text{0.2}}$
        & $\text{97.8}_{\pm\text{0.4}}$
        & $\text{97.6}_{\pm\text{0.3}}$
        & $\text{96.6}_{\pm\text{0.5}}$
        & 97.6 \\

        \textbf{\ours-Decoupled}
        & $\text{99.0}_{\pm\text{0.3}}$
        & $\text{98.0}_{\pm\text{0.2}}$
        & $\text{97.8}_{\pm\text{0.5}}$
        & $\text{97.0}_{\pm\text{0.4}}$
        & 98.0 \\
        \bottomrule
    \end{tabular}
    % \vspace{-0.1in}
\end{table}


\begin{table}[!t]
    \centering
    \caption{
    Closed-loop success (\%) on real-world bimanual manipulation tasks.
    }
    % \vspace{-0.1in}
    \label{tab:real-world-native}
    \small
    \begin{tabular}{@{}lcccc@{}}
        \toprule
        \textbf{Method}
        & \textbf{Toast}
        & \textbf{Cube}
        & \textbf{Cups}
        & \textbf{Mean} \\
        \midrule
        \ours-VTA
        & 92.0 & 90.0 & 94.4 & 92.1 \\
        \ours-Joint
        & 90.0 & 94.0 & 91.7 & 91.9 \\
        \bottomrule
    \end{tabular}
    \vspace{-0.2in}
\end{table}

\begin{table*}[t]
    \centering
    \caption{
    Frozen inverse dynamics transfer to four LIBERO-90 targets.
    Target rows use task-tuned video predictors while source-trained IDMs
    remain frozen. Mean is averaged over the four target tasks.
    }
    % \vspace{-0.1in}
    \label{tab:component-transfer}
    \small
    \setlength{\tabcolsep}{5pt}
    \begin{tabular}{@{}llccccc@{}}
        \toprule
        \textbf{Video / policy}
        & \textbf{Action component (frozen)}
        & \shortstack{\textbf{64}\\Composition}
        & \shortstack{\textbf{74}\\Retarget}
        & \shortstack{\textbf{21}\\Object/grasp}
        & \shortstack{\textbf{45}\\Scene shift}
        & \shortstack{\textbf{Transfer}\\\textbf{mean}} \\
        \midrule

        \multicolumn{2}{l}{
            \textbf{Task-tuned VTA (finetuned policy reference)}
        }
        & 94\%
        & 86\%
        & 92\%
        & 100\%
        & 93.0\% \\

        \midrule

        Task-tuned video
        & Full-context IDM, demo-only
        & 88\%
        & 90\%
        & 0\%
        & 0\%
        & 44.5\% \\

        Task-tuned video
        & Full-context IDM, mixed
        & 54\%
        & 92\%
        & 4\%
        & 38\%
        & 47.0\% \\

        \midrule

        Task-tuned video
        & Local-context IDM, demo-only
        & 36\%
        & 50\%
        & 0\%
        & 0\%
        & 21.5\% \\

        Task-tuned video
        & Local-context IDM, CF-only
        & 80\%
        & 92\%
        & 90\%
        & 76\%
        & \textbf{84.5\%} \\

        Task-tuned video
        & Local-context IDM, mixed
        & 74\%
        & 86\%
        & 88\%
        & 88\%
        & \textbf{84.0\%} \\

        \midrule

        \multicolumn{2}{l}{
            \textbf{Original VTA (unadapted policy reference)}
        }
        & 42\%
        & 0\%
        & 0\%
        & 0\%
        & 10.5\% \\

        \bottomrule
    \end{tabular}
    \vspace{-0.1in}
\end{table*}


% ---------------------------------------------------------
% Source-task composition on LIBERO-Long
% ---------------------------------------------------------
\begin{table}[!t]
    \centering
    \caption{
    Source-task composition on LIBERO-Long, all action components evaluated with LIBERO-Long video
    predictor.
    }
    % \vspace{-0.1in}
    \label{tab:component-source}
    \footnotesize
    \begin{tabular}{@{}llc@{}}
        \toprule
        \textbf{Action component}
        & \textbf{Supervision}
        & \textbf{Success (\%)} \\
        \midrule

        Full-context IDM
        & Demo-only (Native VTA)
        & 97.8 \\

        Full-context IDM
        & Mixed
        & 92.2 \\

        \midrule

        Local-context IDM
        & Demo-only
        & 25.8 \\

        Local-context IDM
        & CF-only
        & 90.6 \\

        Local-context IDM
        & Mixed
        & \textbf{94.4} \\

        \bottomrule
    \end{tabular}
    % \vspace{-0.1in}
\end{table}



\begin{table}[t]
    \centering
    \caption{
    Local-context FDM prediction on counterfactual transitions.
    MSE values are reported in units of $10^{-3}$.
    Outcome accuracy measures identification of the matching realized future
    by RGB MSE among $K$ same-state outcomes.
    }
    \label{tab:fdm}
    \small
    \begin{tabular}{@{}lcccc@{}}
        \toprule
        \textbf{Supervision}
        & \textbf{RGB MSE $\downarrow$}
        & \textbf{Acc.} $K{=}2$ $\uparrow$
        & \textbf{Acc.} $K{=}16$ $\uparrow$ \\
        \midrule
        Demo-only
        & 14.35
        & 68.3
        & 21.1 \\
        CF-only
        & \textbf{9.40}
        & \textbf{93.6}
        & \textbf{71.3} \\
        Mixed
        & 9.62
        & 91.8
        & 67.7 \\
        \bottomrule
    \end{tabular}
    \vspace{-0.2in}
\end{table}

\subsection{Native Policy Performance}
\label{sec:native-policy-performance}

\paragraph{Broad LIBERO evaluation.}
% We first test whether the shared foundation supports strong native policies
% across different generation orders, and report the average success rate over three random seeds in Table~\ref{tab:libero-broad}.
We compare the four interaction programs across the LIBERO suites and
report the average success rate over three random seeds in
Table~\ref{tab:libero-broad}.
The principal \ours configurations achieve strong performance across all
four suites, establishing that the shared foundation supports multiple
world--action generation orders.

\paragraph{Backbone initialization ablation.}
Holding the VTA architecture, task supervision, and downstream training
procedure fixed, the original Wan2.2-5B initialization reaches 68.4\%
success on LIBERO-Long, while our robot-video-pretrained causal backbone
reaches 97.8\% (+29.4 points). This comparison measures the combined effect
of robot-video pretraining and causal adaptation, rather than attributing the
gain to either factor individually.


\paragraph{Real-world native policies.}



% The physical evaluation tests the same framework under multi-view bimanual,
% deformable, contact-rich, and multi-object manipulation.
Table~\ref{tab:real-world-native} reports closed-loop success for VTA and
Joint on the three bimanual tasks of Table~\ref{tab:real-world-tasks}, which
test the same framework under multi-view, deformable, contact-rich,
and multi-object manipulation.


\subsection{Component Composability}
\label{sec:component-composability}


% ---------------------------------------------------------
% Target-task transfer on LIBERO-90
% ---------------------------------------------------------


We next study whether action-grounding components remain useful when the visual predictor changes. Full-context and local-context IDMs are evaluated under transition-supervision regimes, separating the effects of broader local-transition data from the restriction of task-level context.

\paragraph{Frozen action-component reuse.}
For each LIBERO-90 target, the visual producer is taken from a task-tuned
VTA model while source-trained action components remain frozen. Thus, the
experiment measures reuse of an action component after visual-task
adaptation, rather than target-task learning without action supervision.
All reused executors are evaluated with the same target predictor and
rollout protocol. Table~\ref{tab:component-source} reports source-task
composition on LIBERO-Long.

Conditioning scope and supervision interact strongly
(Tables~\ref{tab:component-source} and~\ref{tab:component-transfer}).
Mixed supervision raises mean target success from 21.5\% to 84.0\%
for local-context IDMs, compared with 44.5\% to 47.0\% for
full-context IDMs. Under matched mixed supervision, local conditioning
improves target success by 37.0 percentage points while retaining
94.4\% source-task success.
% Neither local conditioning nor broader
% supervision alone matches their combination in this evaluation.
Counterfactual-only and mixed supervision both substantially improve
local-context IDM transfer, achieving comparable mean success of 84.5\%
and 84.0\%, respectively, across the four LIBERO-90 targets. These
results show that broader transition coverage supports the reuse of
local-context dynamics components after visual task adaptation.

The per-task results also reveal uneven transfer. The demonstration-only
full-context IDM performs well on Tasks 64 and 74 but obtains 0\%
success on Tasks 21 and 45, whereas CF-only and mixed local-context
IDMs remain effective across all four targets. This contrast highlights
their more consistent reuse across the evaluated task changes.


\subsection{Action-Conditioned Forward Dynamics}

We evaluate local-context FDMs on 2,560 counterfactual futures from
160 contexts across ten LIBERO tasks, with 16 action branches per context.
Models share the same initial observations, actions, and realized futures.
In addition to RGB prediction error, we measure whether the predicted future
preserves the identity of its conditioning action through outcome
identification at $K=2$ and $K=16$.

As shown in Table~\ref{tab:fdm}, CF-only and mixed supervision reduce RGB MSE
by 34.5\% and 33.0\%, respectively, relative to demonstration-only training.
CF-only supervision raises outcome accuracy from 68.3\% to 93.6\% for
$K=2$, and from 21.1\% to 71.3\% for $K=16$. Mixed supervision shows a
similar improvement, reaching 91.8\% and 67.7\%, respectively. These results
show that broader transition supervision improves both visual prediction
accuracy and the ability to distinguish action-dependent futures.